\chapter{Finance, product control et fiscalité}\label{ch:fm:finance-product-control-and-tax}

Entre mi-2012 et mi-2015, l'administration fiscale danoise a versé des remboursements d'impôt sur les dividendes à 4\,170 demandeurs, pour un peu moins de 12.1~milliards de couronnes (environ
1.4~milliard de livres), à des demandeurs dont les transactions sur actions avaient été organisées selon les modalités du cum-ex. Aucune de ces demandes n'était valable au regard du droit danois. L'organisateur
principal du système figure parmi les personnes condamnées au Danemark~; un jury new-yorkais a accordé à l'administration 500~millions de dollars contre certaines des parties. Mais en
octobre 2025, la Commercial Court de Londres a rejeté les demandes civiles contestées de l'administration contre les défendeurs établis sur place, au motif qu'elle n'avait pas démontré avoir été
trompée~: ses contrôles de paiement des demandes de remboursement, a estimé le juge, «~étaient si peu rigoureux qu'ils étaient presque inexistants~». Une règle fiscale détermine quelles opérations
existent, et une transaction dont le profit est un remboursement d'impôt pose une question juridique avant d'être une question de trading. Un chiffre de P\&L relève du même raisonnement~: quelqu'un doit le valider,
et ce qu'il vérifie en décide la signification.

\section{Le P\&L quotidien~: flash, définitif et validation}

Chaque journée de trading se termine par deux chiffres. Le premier arrive quelques minutes après la clôture, à partir des systèmes et des valorisations propres du desk~; le second, le
lendemain matin, une fois toutes les transactions saisies, les frais prélevés et les valorisations vérifiées par une personne autre que le trader (le product control du livre~6).

\begin{definition}[P\&L flash, validation du P\&L]\label{def:fm:finance-product-control-and-tax:flash}
Le \emph{P\&L flash}\index{P\&L flash} est une estimation du P\&L de la journée établie par le desk peu après la clôture, à partir des transactions saisies avant une
heure de coupure et des valorisations de fin de journée propres au desk. La \emph{validation du P\&L}\index{validation du P\&L} est l'approbation quotidienne du P\&L définitif
par le product control et le responsable du desk, après vérification indépendante des prix, frais et ajustements, chaque exception étant expliquée et consignée.
\end{definition}

\begin{definition}[Passage du flash au définitif]\label{def:fm:finance-product-control-and-tax:walk}
Le \emph{passage du flash au définitif}\index{passage du flash au définitif} répartit l'écart entre le P\&L définitif et le \omterm{def:fm:finance-product-control-and-tax:flash}{P\&L flash} en catégories, dont chacune a un
responsable~: les mouvements de marché après les valorisations du flash, les transactions saisies après l'heure de coupure, les frais et commissions, et les ajustements de valorisation issus de la
vérification indépendante des prix et des réserves. Ce que les catégories n'expliquent pas est signalé comme inexpliqué.
\end{definition}

Construit comme une suite de réévaluations d'un même livre, le passage se boucle par construction~: on fait passer les valorisations du flash à la clôture sur les transactions du flash, on ajoute
les transactions tardives à la clôture, on impute les frais, puis on fait passer les valorisations de la clôture aux valorisations vérifiées. Chaque étape est la différence de deux totaux d'un registre firm.pnl
(livre~1), et les étapes se compensent de proche en proche, du flash au définitif (le \cref{lst:fm:finance-product-control-and-tax:walk}). Un passage qui ne se boucle pas
est lui-même une exception~: une transaction, une valorisation ou des frais figure dans un chiffre et pas dans l'autre.

\omcode{code/firm/signoff/firm_signoff.py}{46}{73}{Le passage du flash au définitif comme suite de réévaluations des registres firm.pnl, et les règles qui bloquent une
validation.}\label{lst:fm:finance-product-control-and-tax:walk}

\section{Tutoriel~: une journée du flash au définitif}\label{tut:fm:finance-product-control-and-tax:tut}

\begin{tutorial}
\textbf{Objectif.} Produire le \omterm{def:fm:finance-product-control-and-tax:flash}{P\&L flash} et le P\&L définitif d'une journée, passer de l'un à l'autre et décider si la journée peut être validée. \textbf{État final~:}
le \cref{tab:fm:finance-product-control-and-tax:walk} et le graphique du passage (la \cref{fig:fm:finance-product-control-and-tax:walk}).
\begin{tutsteps}
\item \textbf{Le portefeuille.} Dix actions détenues en début de journée, pour 66.5~millions de dollars en valeur brute, huit en position longue et deux en position courte~; trois sont peu liquides
(données synthétiques~; \texttt{fm\_finance.day}).
\item \textbf{Les transactions.} Quarante exécutions au cours de la journée~; les six postérieures à l'heure de coupure de 16:15 ne figurent pas dans le flash. Chaque exécution donne lieu à des frais de 5~points de base.
\item \textbf{Les valorisations.} Le flash utilise les valorisations de 16:00 du desk~; les valorisations de clôture suivent~; la vérification indépendante des prix (IPV) ramène les valorisations des trois positions peu liquides
au bord d'une bande de 0.5~\% autour d'un consensus (\texttt{firm.pnlexplain.ipv}, livre~6).
\item \textbf{Les règles.} Suspendre la validation si le définitif s'écarte du flash de plus de 100\,000~dollars ou de 10~\% du flash, selon le plus élevé des deux~;
si les ajustements de valorisation dépassent 100\,000~dollars, ou si les transactions tardives dépassent 50\,000~dollars~; ou si quoi que ce soit reste inexpliqué.
\end{tutsteps}
\end{tutorial}

\begin{table}[ht]
\centering\small
\begin{tabular}{lr}
\toprule
étape & \$ \\
\midrule
\omterm{def:fm:finance-product-control-and-tax:flash}{P\&L flash} & 475\,100 \\
mouvements de marché après le flash & +63\,910 \\
transactions tardives (6 exécutions) & +1\,580 \\
frais & $-$11\,670 \\
ajustements de valorisation & $-$149\,787 \\
\midrule
P\&L définitif & 379\,133 \\
inexpliqué & 0 \\
\bottomrule
\end{tabular}
\caption{Le \omterm{def:fm:finance-product-control-and-tax:walk}{passage du flash au définitif} de la journée. Données~: \texttt{fm\_finance.the\_walk}.}\label{tab:fm:finance-product-control-and-tax:walk}
\end{table}

\begin{omfigure}
\begin{tikzpicture}
\begin{axis}[width=10.5cm,height=5.4cm,ybar stacked,bar width=20pt,ylabel={\small milliers de \$},ymin=0,ymax=600,xtick={0,1,2,3,4,5},
  xticklabels={flash,{marché},{tardives},frais,{valorisation},définitif},xmin=-0.6,xmax=5.6,
  tick label style={font=\footnotesize},grid=major,grid style={gray!20},table/col sep=comma]
\addplot[draw=none,fill=none] table[x=pos,y=base]{figdata/desk/15-finance-product-control-and-tax/walk.csv};
\addplot[fill=omDef!55,draw=omDef] table[x=pos,y=height]{figdata/desk/15-finance-product-control-and-tax/walk.csv};
\end{axis}
\end{tikzpicture}

\omcaption{Le passage d'un \omterm{def:fm:finance-product-control-and-tax:flash}{P\&L flash} de 475\,100~dollars à un P\&L définitif de 379\,133~dollars~: mouvements de marché après les valorisations du flash, six transactions tardives, frais et
ajustements de valorisation sur trois positions peu liquides. Chaque barre couvre l'étape qu'elle explique. Données~: \texttt{fm\_finance.the\_walk}.}\label{fig:fm:finance-product-control-and-tax:walk}
\end{omfigure}

Le P\&L définitif est inférieur de 20.2~\% au flash, soit 95\,967~dollars, ce qui reste dans la tolérance globale de 100\,000~dollars~; une règle portant sur le seul total la validerait. La
règle par catégorie l'arrête~: les ajustements de valorisation de 149\,787~dollars dépassent leur tolérance de 100\,000~dollars, si bien que la journée est remontée avec le motif consigné
(\texttt{fm\_finance.day\_signoff}). L'ajustement a une cause claire. Le desk a clôturé sa plus grande position peu liquide, 73\,000 actions, à
123.49~dollars, contre un consensus de 122.01~dollars, soit 1.2~\% de plus~; la vérification a ramené la valorisation au bord de la bande de 0.5~\%, soit 122.63~dollars, et a fait de même, en
de moindre ampleur, sur les deux autres. Un desk dont les valorisations de clôture sur des positions peu liquides se tiennent au bord favorable du consensus jour après jour est le
schéma que le product control est précisément là pour repérer.

\begin{method}[Validation du P\&L quotidien]\label{met:fm:finance-product-control-and-tax:sign}
\begin{enumerate}
\item Publier le flash avec son heure de coupure et la source de ses valorisations.
\item Enregistrer les transactions tardives, les frais et les ajustements~; vérifier les valorisations de manière indépendante et appliquer les réserves (évaluation prudente du livre~6).
\item Établir le \omterm{def:fm:finance-product-control-and-tax:walk}{passage du flash au définitif} par catégorie~; le passage doit se boucler.
\item Appliquer les seuils sur le total et sur chaque catégorie~; remonter chaque exception avec un motif écrit.
\item Valider, et conserver la piste d'audit~: qui a validé quoi, quand et sur quels éléments. La piste d'audit ne s'alimente que par ajout.
\end{enumerate}
\end{method}

\section{Contrôle de la valorisation dans une société de trading}

L'autorité du product control repose sur la même indépendance que celle de la fonction risque (chapitre 12)~: il rend compte en dehors des desks, il contrôle les valorisations
qui entrent dans les livres et registres, et ses ajustements ne peuvent être annulés par le trader. Dans une banque, la fonction est large et réglementée~; dans une
société pour compte propre ou un fonds, elle peut tenir à deux personnes ou se confier à un administrateur (chapitre 4), et les tentations sont les mêmes. Les valorisations d'un trader sur
des positions peu liquides sont une estimation de la personne dont la rémunération en dépend (chapitre 10). Les contrôles sont ceux du livre~6~: vérification indépendante face au
consensus ou aux transactions observées, bandes de tolérance, réserves pour ce qui ne peut être vérifié, et P\&L inexpliqué examiné plutôt qu'absorbé.

\begin{remark}[Le flash comme contrôle]\label{rem:fm:finance-product-control-and-tax:flash}
Un flash toujours au-dessus du définitif est une information~: il signifie que les valorisations propres du desk sont optimistes, que la saisie de ses transactions est tardive, ou que ses frais sont
oubliés. Suivre les catégories du passage sur plusieurs mois, et non seulement les exceptions du jour, fait de la validation une mesure de la façon dont un desk
connaît sa propre position.
\end{remark}

\section{Les impôts qui façonnent les stratégies}

Les impôts entrent dans le rendement d'une stratégie sous forme de coûts sur ses transactions, sur ses revenus et sur ses gains. Les deux premiers sont payés quel que soit le résultat et croissent avec
l'activité~; ils décident des stratégies qui existent sur chaque marché.

\begin{definition}[Frein fiscal]\label{def:fm:finance-product-control-and-tax:drag}
Le \emph{frein fiscal}\index{frein fiscal} (tax drag) d'une stratégie est la réduction de son rendement annuel causée par les impôts~: taxes sur les transactions portant sur son volume de rotation,
retenues à la source sur les dividendes et intérêts qu'elle perçoit et qu'elle ne peut pas récupérer, et impôts sur ses gains, pour un investisseur et une juridiction donnés.
\end{definition}

\begin{proposition}[Le volume de rotation que l'impôt absorbe]\label{prop:fm:finance-product-control-and-tax:half}
Une stratégie dont le rendement net est $r$ par an avant une taxe sur les transactions, qui applique un taux $t$ sur ses achats, lesquels représentent $T$ fois son capital par an, obtient
$r-Tt$ après l'impôt. L'impôt absorbe la moitié de son rendement net pour
\[
T_{1/2}=\frac{r}{2t},
\]
et la totalité à deux fois ce volume de rotation.
\end{proposition}

\begin{proof}
L'impôt payé dans l'année est le taux multiplié par les achats, soit $tT$ par unité de capital~; poser $tT=r/2$ donne $T_{1/2}$.
\end{proof}

Pour une stratégie qui rapporte 8~\% par an, une taxe de 10 points de base absorbe la moitié de son rendement à un volume de rotation de 40 fois par an, 20 points de base à 20 fois, et
50 points de base, le taux britannique sur les achats d'actions (\cref{dat:fm:finance-product-control-and-tax:tax}), à 8 fois
(\cref{fig:fm:finance-product-control-and-tax:turnover}). Un teneur de marché ou un arbitragiste statistique qui fait tourner son livre quarante fois par an ne peut pas
négocier un instrument taxé à 50 points de base~; il négocie là où la taxe ne s'applique pas, là où la réglementation le permet, ou ne négocie pas du tout.

\begin{omfigure}
\begin{tikzpicture}
\begin{axis}[width=10.5cm,height=5.2cm,xlabel={\small volume de rotation (achats par an sur capital)},ylabel={\small rendement après l'impôt (\%)},xmin=0,xmax=50,
  ymin=-10,ymax=9,tick label style={font=\footnotesize},grid=major,grid style={gray!20},table/col sep=comma,
  legend cell align=left,legend style={font=\scriptsize,at={(0.5,-0.3)},anchor=north,/tikz/every even column/.append style={column sep=8pt}},legend columns=3]
\addplot[omThm,thick] table[x=turnover,y=t10]{figdata/desk/15-finance-product-control-and-tax/turnover.csv};
\addplot[omProp,thick,dashed] table[x=turnover,y=t20]{figdata/desk/15-finance-product-control-and-tax/turnover.csv};
\addplot[omDef,very thick,dotted] table[x=turnover,y=t50]{figdata/desk/15-finance-product-control-and-tax/turnover.csv};
\draw[gray,dotted] (axis cs:0,4) -- (axis cs:50,4);
\draw[black,thin] (axis cs:0,0) -- (axis cs:50,0);
\legend{10 pb,20 pb,50 pb}
\end{axis}
\end{tikzpicture}

\omcaption{Le rendement, après une taxe sur les achats, d'une stratégie qui rapporte 8~\% par an avant celle-ci, en fonction du volume de rotation, pour trois taux d'impôt. La
ligne pointillée représente la moitié du rendement~: atteinte à des volumes de 40, 20 et 8 (\cref{prop:fm:finance-product-control-and-tax:half}). Données~:
\texttt{firm.signoff.tax\_drag}.}\label{fig:fm:finance-product-control-and-tax:turnover}
\end{omfigure}

La retenue à la source frappe l'autre extrémité~: les stratégies de revenu. Une stratégie value avec un rendement de dividende de 3.5~\% perd 0.53~\% par an à
une retenue de 15~\% non récupérable, et un volume de rotation de 0.8 par an ne lui coûte que 0.40~\% à 50 points de base~; pour un livre d'arbitrage statistique qui
fait tourner son livre quarante fois, les mêmes 50 points de base coûtent 20~\% par an (\cref{tab:fm:finance-product-control-and-tax:tax}). Les taux du tableau sont
des données illustratives~; l'encadré daté donne les règles telles qu'elles sont en vigueur.

\begin{table}[ht]
\centering\small
\begin{tabular}{lrrrrrr}
\toprule
 & rendement & volume de & rendement de & régime 1 & régime 2 & régime 3 \\
stratégie & brut & rotation & dividende & & & \\
\midrule
arbitrage statistique & 12.0 & 40 & 1.0 & 11.85 & $-$8.00 & 3.85 \\
momentum & 9.0 & 6 & 1.5 & 8.77 & 6.00 & 7.58 \\
value & 7.0 & 0.8 & 3.5 & 6.48 & 6.60 & 6.31 \\
\bottomrule
\end{tabular}
\caption{Rendements après impôt de trois stratégies stylisées sous trois régimes illustratifs (\%)~: (1) aucune taxe sur les transactions et retenue de 15~\%~;
(2) 0.5~\% sur les achats~; (3) 0.2~\% sur les achats et retenue de 15~\%. Le volume de rotation est le montant des achats par an sur capital. Données~: \texttt{fm\_finance.tax\_table}.}\label{tab:fm:finance-product-control-and-tax:tax}
\end{table}

\begin{definition}[Opérations cum-ex, règle des ventes fictives]\label{def:fm:finance-product-control-and-tax:cumex}
Les \emph{opérations cum-ex}\index{opérations cum-ex} consistent à négocier des actions autour d'une date de dividende, de sorte qu'une transaction conclue à la date de déclaration du dividende ou avant
soit réglée après la date d'enregistrement, ce qui sert à générer des demandes de remboursement de l'impôt sur le dividende par plus d'une partie pour un même dividende,
ou par des parties qui ne l'ont jamais reçu. La \emph{règle des ventes fictives}\index{règle des ventes fictives} (wash-sale rule) interdit la déduction d'une perte fiscale sur la cession d'un titre lorsque le vendeur
acquiert le même titre, ou un titre substantiellement identique, dans une période déterminée autour de la cession.
\end{definition}

Le cas danois est la première forme de transaction poussée à son terme~: une transaction dont le contenu économique était un remboursement d'impôt, réclamé sur des actions que les demandeurs
n'avaient pas détenues. La cour britannique a jugé que chacune des 4\,170 demandes était invalide au regard du droit danois, et a estimé que les concepteurs du dispositif n'avaient aucun fondement raisonnable pour
les croire valables~; elle a néanmoins rejeté les demandes civiles contestées pour tromperie, au motif que les procédures de l'administration ne reposaient pas sur
ce qui lui avait été dit. Le livre~1 décrit les affaires allemandes, qui se sont terminées par des condamnations pénales. Pour une société de trading, la leçon est celle de l'accroche du chapitre~: une stratégie
dont le rendement dépend d'une position fiscale est approuvée par les fonctions fiscale et juridique de la société avant d'être négociée, et son profit est audité comme une
position fiscale, et non comme une position de trading.

\begin{dated}{2026-09}{Règles fiscales qui façonnent le trading}\label{dat:fm:finance-product-control-and-tax:tax}
\textbf{États-Unis}~: en vertu du 26 U.S.C. 1256, les contrats à terme réglementés et les autres contrats de la section 1256 détenus en fin d'année sont considérés comme vendus à leur juste
valeur de marché le dernier jour ouvré, et leurs gains et pertes sont à 60~\% à long terme et à 40~\% à court terme, quelle que soit la durée de détention~; en vertu du 26
U.S.C. 1091, une perte sur la cession d'actions ou de titres est refusée si des titres substantiellement identiques sont acquis dans les 30 jours qui précèdent ou
suivent la cession, sauf si le vendeur est un dealer agissant dans le cours normal de son activité. \textbf{Royaume-Uni}~: un impôt ou un droit de 0.5~\% est
généralement exigible sur les achats d'actions (Stamp Duty Reserve Tax, ou droit de timbre sur les transferts sur papier). Les taux sur les gains dépendent de l'investisseur et
ne sont pas précisés ici. Cet encadré résume des textes publics et n'est pas un conseil fiscal.
\end{dated}

Avec des taux illustratifs de 37~\% à court terme et de 20~\% à long terme, la règle 60/40 impose aux gains sur contrats à terme un taux mixte de 26.8~\%, quelle que soit la
période de détention, ce qui explique pourquoi le même rendement de suivi de tendance peut valoir davantage après impôt sur contrats à terme que sur des actions détenues pendant des semaines
(\texttt{firm.signoff.blended}).

\section{Entités juridiques et prix de transfert}

Une société qui négocie dans plusieurs pays le fait par l'intermédiaire de plusieurs entités juridiques~: une entité de trading où siègent ses traders, une entité qui détient
les adhésions aux bourses, une société de gestion qui emploie les chercheurs, un fonds ou un véhicule pour compte propre qui détient le capital. Chacune est imposée
là où elle est résidente, et chacune facture aux autres ce qu'elle fournit~: la recherche, l'exécution, la technologie, le capital.

\begin{definition}[Prix de transfert]\label{def:fm:finance-product-control-and-tax:tp}
Le \emph{prix de transfert}\index{prix de transfert} (transfer pricing) est la fixation des prix des transactions entre entités d'un même groupe (services,
concessions de propriété intellectuelle, financement, partage des profits de trading), que les administrations fiscales exigent qu'il soit celui que
accepteraient des parties indépendantes dans des circonstances comparables~: le principe de pleine concurrence.
\end{definition}

Pour une société de trading, les cas difficiles sont ceux que ce livre a déjà rencontrés~: qui détient la stratégie (chapitre 11), où sont prises les décisions qui génèrent
le profit, et ce qu'un desk d'un pays paie à une équipe de recherche située dans un autre. Les Transfer Pricing Guidelines de l'OCDE (édition 2022) sont le
texte de référence sur le principe de pleine concurrence~; la documentation qu'elles exigent est tenue par la finance, et c'est la raison pour laquelle l'organigramme d'une société et son
schéma des entités juridiques sont dessinés séparément et doivent concorder.

\begin{omfigure}
\begin{tikzpicture}[font=\footnotesize,>=stealth,box/.style={draw,rounded corners,align=center,minimum height=0.75cm,text width=2.8cm}]
\node[box,fill=gray!15] (h) at (0,2.2) {société holding};
\node[box,fill=omDef!15] (t) at (-4.3,0) {entité de trading\\{\scriptsize traders, capital}};
\node[box,fill=omProp!15] (r) at (0,0) {société de recherche\\{\scriptsize chercheurs, code}};
\node[box,fill=omThm!15] (m) at (4.3,0) {entité membre\\{\scriptsize accès à la bourse}};
\draw[->] (h) -- (t); \draw[->] (h) -- (r); \draw[->] (h) -- (m);
\draw[->,omProp,thick] (r) -- (t); \node[font=\scriptsize] at (-2.15,-0.7) {frais de recherche};
\draw[->,omThm,thick] (m.south) to[bend left=25] node[below,font=\scriptsize,text=black] {frais d'exécution} (t.south);
\end{tikzpicture}

\omcaption{Un groupe stylisé d'entités et les facturations intragroupe que le \omterm{def:fm:finance-product-control-and-tax:tp}{prix de transfert} doit fixer selon le principe de pleine concurrence~: l'entité de trading paie la
société de recherche pour ses modèles et l'entité membre pour son accès. Schématique.}\label{fig:fm:finance-product-control-and-tax:entities}
\end{omfigure}

\section{Construction~: le moteur de validation}\label{bld:fm:finance-product-control-and-tax:signoff}

\begin{build}
\bfield{Rôle} Le P\&L quotidien du flash au définitif, le passage entre les deux, les règles de validation avec piste d'audit, et un calculateur de \omterm{def:fm:finance-product-control-and-tax:drag}{frein fiscal} dont
les taux sont des entrées.
\bfield{Interface} \texttt{firm.signoff}~: \texttt{day\_book} (sur \texttt{firm.pnl}), \texttt{verified\_marks} (sur \texttt{firm.pnlexplain.ipv}),
\texttt{walk}, \texttt{exceptions}, \texttt{Trail}, \texttt{sign}~; \texttt{tax\_drag}, \texttt{half\_turnover}, \texttt{blended}.
\bfield{Règles} Le passage s'additionne exactement, en unités comptables~; une validation comportant des exceptions est remontée, jamais signée~; la piste d'audit est en ajout seul~; les taux
fiscaux sont toujours des entrées.
\bfield{Tests d'acceptation} \texttt{code/firm/signoff/tests/}~: un passage à deux instruments contrôlé ligne à ligne~; valorisations vérifiées à la limite de la bande~; les
règles d'exception~; les formules fiscales.
\bfield{Pour aller plus loin} Réserves pour le spread bid-offer et l'incertitude de modèle~; un mois de passages avec les tendances par catégorie et par desk~; les demandes de remboursement de retenue à la source avec un décalage.
\end{build}

\begin{omsources}
\item Skatteforvaltningen v Solo Capital Partners LLP and others [2025] EWHC 2364 (Comm), jugement du 2 octobre 2025.
\item 26 U.S.C. 1256 et 1091~; GOV.UK, «~Tax when you buy shares~».
\item OECD, \emph{Transfer Pricing Guidelines for Multinational Enterprises and Tax Administrations 2022}.
\end{omsources}

\section{Exercices}

\begin{exercise}[$\star$]\label{exo:fm:finance-product-control-and-tax:1}
Vérifier que les étapes du \cref{tab:fm:finance-product-control-and-tax:walk} s'additionnent bien jusqu'au définitif.
\end{exercise}

\begin{exercise}[$\star$]\label{exo:fm:finance-product-control-and-tax:2}
Pour quelle rotation annuelle une taxe de 20 points de base retire-t-elle la moitié du rendement d'une stratégie qui rapporte 8~\% par an~? Et la totalité~?
\end{exercise}

\begin{exercise}[$\star$]\label{exo:fm:finance-product-control-and-tax:3}
Avec la règle 60/40 et des taux illustratifs de 37~\% et 20~\%, quel est le taux mixte sur un gain sur contrats à terme~?
\end{exercise}

\begin{exercise}[$\star\star$]\label{exo:fm:finance-product-control-and-tax:4}
Pourquoi la journée du tutoriel a-t-elle passé la règle sur le total et échoué à la règle sur les ajustements de valorisation~? Quelle règle est la bonne~?
\end{exercise}

\begin{exercise}[$\star\star$]\label{exo:fm:finance-product-control-and-tax:5}
Un trader vend une action à perte le 20 décembre et la rachète le 5 janvier. Que fait la \omterm{def:fm:finance-product-control-and-tax:cumex}{règle des ventes fictives} (wash-sale rule), et qu'en est-il si le trader est un dealer
agissant dans le cadre normal de son activité~?
\end{exercise}

\begin{exercise}[$\star\star$]\label{exo:fm:finance-product-control-and-tax:6}
Pourquoi la cour britannique a-t-elle rejeté les prétentions de l'administration fiscale alors qu'elle n'a jugé valable aucune des demandes de remboursement~?
\end{exercise}

\begin{exercise}[$\star\star\star$]\label{exo:fm:finance-product-control-and-tax:7}
\emph{Programmation.} Lancer \texttt{fm\_finance.the\_walk(17)}. Quels sont le flash et le définitif, et quelle règle arrête la validation~?
\end{exercise}

\begin{exercise}[$\star\star\star$]\label{exo:fm:finance-product-control-and-tax:8}
\emph{Trouver la faille.} «~Notre flash est toujours à moins de 5~\% du définitif, donc nos valorisations sont bonnes.~»
\end{exercise}

\section{Problème~: La rotation que la taxe prélève}

\begin{problem}[{Problème du week-end --- la rotation que la taxe prélève}]\label{pb:fm:finance-product-control-and-tax:1}
Un nouveau directeur financier doit fixer les règles de validation de la société et vérifier si ses stratégies résistent aux taxes des marchés sur lesquels elle prévoit d'entrer.

\medskip\noindent\textbf{Partie I --- Le P\&L quotidien.}
\begin{enumerate}
\item Définir le \omterm{def:fm:finance-product-control-and-tax:flash}{P\&L flash}, la \omterm{def:fm:finance-product-control-and-tax:flash}{validation du P\&L} et le \omterm{def:fm:finance-product-control-and-tax:walk}{passage du flash au définitif}.
\item Pourquoi un passage construit à partir des réévaluations s'additionne-t-il exactement~?
\item Donner le passage de la journée du tutoriel.
\item Quelles règles arrêtent sa validation, et pourquoi~?
\item Expliquer l'ajustement de valorisation sur la position illiquide la plus importante.
\end{enumerate}

\medskip\noindent\textbf{Partie II --- Contrôle.}
\begin{enumerate}[resume]
\item Qu'est-ce qui rend le product control indépendant~?
\item Que révèle un flash qui est toujours supérieur au définitif~?
\item Que doit enregistrer la piste d'audit de la validation~?
\end{enumerate}

\medskip\noindent\textbf{Partie III --- Fiscalité.}
\begin{enumerate}[resume]
\item Définir le \omterm{def:fm:finance-product-control-and-tax:drag}{frein fiscal} (tax drag).
\item Énoncer et démontrer la \cref{prop:fm:finance-product-control-and-tax:half}.
\item Donner les rotations annulant la moitié du rendement à 10, 20 et 50 points de base pour une stratégie rapportant 8~\%.
\item Donner les rendements des trois stratégies sous les trois régimes.
\item Énoncer la règle 60/40 et la \omterm{def:fm:finance-product-control-and-tax:cumex}{règle des ventes fictives} (wash-sale rule), ainsi que le taux britannique sur les achats d'actions.
\end{enumerate}

\medskip\noindent\textbf{Partie IV --- Le droit.}
\begin{enumerate}[resume]
\item Définir les \omterm{def:fm:finance-product-control-and-tax:cumex}{opérations cum-ex}.
\item Résumer les demandes de remboursement danoises et le jugement britannique.
\item Que doit faire une société avant de lancer une stratégie dont le rendement dépend d'une position fiscale~?
\item Définir le \omterm{def:fm:finance-product-control-and-tax:tp}{prix de transfert} (transfer pricing) et donner deux facturations intragroupe dans un groupe de trading.
\item Pourquoi l'organigramme et le schéma des entités juridiques doivent-ils concorder~?
\item Énoncer le \emph{résultat nommé}~: la rotation annuelle à laquelle une taxe sur les transactions de $t$ points de base retire la moitié du rendement net d'une stratégie, et
l'écart du flash au définitif qui doit arrêter une validation.
\item En deux phrases, écrire les règles de la société pour la \omterm{def:fm:finance-product-control-and-tax:flash}{validation du P\&L} et pour les stratégies dépendant de la fiscalité.
\end{enumerate}
\end{problem}

\section{Questions d'entretien}

\begin{interviewq}[$\star$ \iqroles{trader}]\label{iq:fm:finance-product-control-and-tax:1}
Quelle est la différence entre le \omterm{def:fm:finance-product-control-and-tax:flash}{P\&L flash} et le P\&L définitif~?
\end{interviewq}

\begin{interviewq}[$\star$ \iqroles{développeur}]\label{iq:fm:finance-product-control-and-tax:2}
Comment vous assureriez-vous que le \omterm{def:fm:finance-product-control-and-tax:walk}{passage du flash au définitif} s'additionne toujours~?
\end{interviewq}

\begin{interviewq}[$\star\star$ \iqroles{chercheur}]\label{iq:fm:finance-product-control-and-tax:3}
Votre stratégie rapporte 6~\% par an avec une rotation de 25. Le marché prélève 20 points de base sur les achats. La mettez-vous en œuvre~?
\end{interviewq}

\begin{interviewq}[$\star\star$ \iqroles{risque}]\label{iq:fm:finance-product-control-and-tax:4}
Les valorisations illiquides d'un desk se situent toujours à la limite favorable de la bande de consensus. Que faites-vous~?
\end{interviewq}

\begin{interviewq}[$\star\star$ \iqroles{trader, chercheur}]\label{iq:fm:finance-product-control-and-tax:5}
Pourquoi le même rendement de suivi de tendance (trend following) peut-il valoir davantage en contrats à terme qu'en actions, pour un contribuable américain~?
\end{interviewq}

\begin{interviewq}[$\star\star\star$ \iqroles{risque, développeur}]\label{iq:fm:finance-product-control-and-tax:6}
Concevez les seuils d'une validation automatisée du P\&L pour cinquante desks.
\end{interviewq}
